\documentclass[11pt]{article}
\usepackage[T1]{fontenc}
\usepackage[utf8]{inputenc}
\usepackage{lmodern}
\usepackage{amsmath,amssymb,amsthm}
\usepackage{mathrsfs}
\usepackage[margin=1in]{geometry}
\usepackage{microtype}
\usepackage[hidelinks]{hyperref}
\newcommand{\MA}{\operatorname{MA}}
\newcommand{\PSH}{\operatorname{PSH}}
\setlength{\emergencystretch}{2em}
\title{P16.3: Mixed-characteristic inputs, model-level progress,\\and unresolved full scope}
\author{Research draft from the frozen assignment and authorized primary sources}
\date{5 October 2026}
\begin{document}
\maketitle
This is a standalone partial research argument for the unchanged original P16.3. It is not a proof of P16.3 in its entire retained domain. The new unconditional conclusion in that domain is a genuine \textbf{model-level} approximation theorem for finite-dual-energy targets. A genuine finite-energy equation theorem is proved below in mixed characteristic \textbf{conditional on the two specific resolution assumptions of FGK}. A concrete intersection example rules out a naive averaging passage through finite covers. None of these conclusions supplies the missing unrestricted model and potential-space inputs. The status is partial and candidate=false.

The historical equal-characteristic surface theorem and its review, supplied in the frozen input, are not rerun or promoted to a general theorem here. The calculus used below is justified by the primary sources in its stated domain. No claim of publication priority or exhaustive contemporary literature coverage is made.

\section*{1. Exact retained problem and normalization}

The original domain throughout is a complete nontrivially discretely valued field \(K\), a smooth geometrically integral projective \(K\)-variety \(X\) of arbitrary dimension \(n\geq 1\), and an actual ample rational line bundle \(L\in\operatorname{Pic}(X)\otimes\mathbb Q\). Put \(R=K^\circ\), \(S=\operatorname{Spec}R\), and

\[
V=(L^n)>0.
\]

A positive multiple \(rL\) is an ample line bundle. A projective embedding by a further multiple and its closure over \(R\) give a relatively nef model reference metric after taking its rational root. Write \(\theta\) for its model curvature. Neither smoothness nor regularity of that closure is asserted or needed for this reference. The reference remains a metric of the actual \(L\); it is not a substitute positivity datum.

Let \(D\) be the real span of model functions on \(X^{\rm an}\), and let

\[
D_\theta=\{v\in D:\theta+dd^cv\text{ is represented by a relatively nef real model class}\}.
\]

Every finite collection of model data has a common projective normal determination. Its mixed intersection measure is genuine: if the special fiber is \(\sum_i b_iE_i\), its coefficient at the corresponding divisorial point \(x_i\) is \(b_i\) times the intersection of the restricted model classes on \(E_i\). On a diagonal it is a positive Radon measure of raw mass \(V\). For real coefficients use multilinearity. Write

\[
\mathcal E(v)=\frac1{n+1}\sum_{j=0}^n
 \int v(\theta+dd^cv)^j\wedge\theta^{n-j},
\qquad e(v)=V^{-1}\mathcal E(v),\qquad m(v)=V^{-1}\MA_\theta(v).
\]

Thus \(m(v)\) is a probability measure and \(e(v+c)=e(v)+c\). For a positive Radon measure \(\mu\), the exact target condition is

\[
\mu(X^{\rm an})=V,\qquad
\mathcal E_\theta^*(\mu)=
\sup_{v\in D_\theta}\left(\mathcal E(v)-\int v\,d\mu\right)<\infty.
\]

Equivalently, \(\nu=\mu/V\) is a probability measure with

\[
A=e^*(\nu)=\sup_{v\in D_\theta}\left(e(v)-\int v\,d\nu\right)<\infty,
\qquad\mathcal E_\theta^*(\mu)=VA.
\]

The entire measure is divided by \(V\). Restrictions to subsets are never renormalized. This condition uses all model tests on the original variety and no skeleton or invariance restriction. Constants cancel only because the mass is exactly \(V\).

For the full original domain, a singular psh cone can be defined by decreasing Fubini--Study nets as in [E], Definition 7.2, excluding the identically negative infinite metric. Stability under max and decreasing limits follows from [E], Lemma 7.4 and Corollary 7.5. Continuous psh metrics agree with uniform limits of nef model metrics by [E], Theorem 7.8. These statements do not establish normalized compactness, a bounded discontinuous local product, or a full finite-energy equation. We do not use a definition of a cone as proof of any of those conclusions.

In residue characteristic zero the published BFJ geometric inputs apply. Removing the restriction additionally requires two distinct regimes: equal characteristic \(\operatorname{char}K=p>0\), and mixed characteristic \(\operatorname{char}K=0,\operatorname{char}\widetilde K=p>0\). A characteristic-p spread-out argument is not an argument about the latter regime.

\section*{2. Exact-source and application ledger}

Labels refer only to the saved versions, not later statements with similar numbering. A binding \texttt{pP:LQ} means PDF page P, line Q in that page's literal cached text; the complete page and its checksum are retained. Recursively cited papers are not represented as independently inspected caches.

\noindent \textbf{[F]} Fang--Gubler--Kunnemann, arXiv:2203.12282v2, identifier 28 February 2024, internal date 29 February 2024. The cache describes the published version as Nagoya Mathematical Journal 259 (2025), 548--561, DOI 10.1017/nmj.2025.8; only the saved arXiv v2 is applied. Section 4 assumes a flat projective integral \textbf{regular} scheme over a complete mixed-characteristic DVR. Definition 4.1 is at p6:L14, Remark 4.2 at p6:L30, Lemma 4.4 at p7:L15, and Theorem 4.7 at p7:L60 continuing on p8. These supply the perturbation-friendly ideals, actual subadditivity, the vertical correction, and generation. The underlying [BMP+24] and [HLS21] results are imported as expressly stated by [F]; their recursive proofs were not read here. The older +-ideal subadditivity conjecture, p5, is not used.

\noindent \textbf{[F], geometric applications.} Theorem 5.1, p8:L59--p10, concerns a regular model on which the generic ample line bundle has an ample extension. Definition 6.1, p10:L18, requires a \textbf{projective} regular domination of every projective model with unchanged generic fiber. Definition 6.3, p10:L34, is the additional embedded resolution assumption; Section 8 uses it to obtain cofinal projective SNC models. Lemma 6.4, p10:L41--p11, and Remark 6.5 provide ample representatives on dominating regular/SNC models, after positive powers. Theorem 7.1, p11:L34, obtains uniform model approximation of every continuous envelope under Definition 6.1. Theorem 8.1, p12:L22, retains \textbf{both} resolution assumptions and SNC-skeleton support of the target. We never attribute arbitrary-support finite-energy surjectivity to that theorem.

\noindent \textbf{[F], important limits and version issue.} Remark 6.2, p10:L23, distinguishes the projective resolution premise from the locally projective morphism cited for arithmetic threefolds. Consequently the paper does not supply an unconditional mixed-characteristic surface instance of Definition 6.1. Section 8, p12:L3--L19, expressly imports BFJ Sections 1--7 and regularization. Its sentence about Sections 4--6 is printed with a [BFJ16] citation, although the equation calculus is in [BFJ15]; we do not use that printed citation as a theorem-number identity. The actual bounded and finite-energy arguments are identified in [B] below and have the field-independent prerequisites checked in Sections 5--8 of this document.

\noindent \textbf{[M]} Boucksom--Jonsson, arXiv:2307.01697v1. Section 1.2.2, p8:L3, supplies the model-intersection synthetic datum for projective Berkovich spaces over a non-Archimedean field, without a residue-characteristic restriction. Section 1.3.2, p11:L18, proves that a model form with ample generic class has positive Bott--Chern class in that datum. It does not assert a general identification of Bott--Chern cohomology with the generic numerical space. Example 2.13 and Definition 2.15, p21, give directed test maxima and test-level orthogonality. Example 2.20, p23:L30, explicitly states that this orthogonality always holds in the projective Berkovich setting, by the relative-volume argument. We import this \textbf{stated conclusion}, including its [BGM22] differentiability estimate; [BGM22] is not a supplied cache and is not independently audited. Theorem 1.33, p17:L31, Proposition 2.2, p18:L39, Theorem 2.23, p23:L42--p24, and Lemma 2.25, p25:L8, apply to precisely this datum. Their conclusion about a completion of measures is not a theorem of surjectivity on unbounded potentials.

\noindent \textbf{[E]} Boucksom--Eriksson, arXiv:1805.01016v3, internal date 11 January 2021. Theorem A, p3:L32, identifies relative volume and energy for \textbf{continuous psh} metrics on an ample line bundle on a geometrically reduced projective scheme over any complete valued field. Definition 7.2, p48:L6, and Theorem 7.8, p49:L12, give the cone and continuous-metric compatibility noted above. Lemma 7.29 and Conjecture 7.30, p54:L3 and L19, explicitly separate envelope continuity from these results. We do not turn Theorem A into an unbounded product theorem.

\noindent \textbf{[S]} Boucksom--Favre--Jonsson, arXiv:1201.0187v2. Its standing residue characteristic is zero. Proposition 5.2, p23:L25, is the ample-representative fact; Definition 7.1, p32:L46--p33, is the full SNC-based psh definition; Theorem 7.10, p34:L35, is normalized compactness; Theorem 8.7, p39:L11, is monotone regularization. Its Sections 1--7 are imported in [F], Section 8. Its characteristic-zero multiplier-ideal argument and vanishing appendix are replaced, not reused.

\noindent \textbf{[B]} Boucksom--Favre--Jonsson, arXiv:1201.0188v1, the actual v1. Theorem 3.1, p16:L31, constructs the bounded mixed operator from model regularization. Proposition 6.2, p26:L28; Definition 6.4 and Lemmas 6.7--6.8, p27; Lemma 6.10, p28:L23; Proposition 6.12, p29:L19; Theorem 7.2 and Corollary 7.3, p30; and Proposition 8.2, p33:L25--p35, supply the energy, truncation, comparison, differentiation, and uniqueness calculus. Their proofs use model intersection integration by parts, Cauchy--Schwarz, regularization, and the indicated compactness/envelope inputs. Theorem A in this v1 has additional algebraizability and support restrictions, and is not invoked as an arbitrary-support theorem. Proposition 4.5 is its sequence-regularization numbering. Lemma 6.10's displayed sign is corrected using its proof: the right side contains the nonpositive potential, not its negative. Proposition 6.12's polynomial degree is at most n+1. Raw uniqueness norms sum to VS, not S. These corrections are explicitly retained.

\noindent \textbf{[G]} Gubler--Jell--Kunnemann--Martin, arXiv:1712.00980v3. Definition 8.1 and Theorem 8.2, p20:L27 and L40, retain equal-characteristic geometric origin over a perfect field and resolution in total dimension d+n. Section 9, p22--p23, additionally retains embedded resolution, gives cofinal SNC models, and explicitly transfers [B], Sections 3--6. Theorem 6.3 is the dimension-three result over perfect fields. Neither it nor the reviewed geometric-origin surface theorem supplies mixed-characteristic projective resolution or unrestricted equal-characteristic geometry. Proposition 2.10, p7:L37, is used below for the pointwise base-ideal envelope argument, which requires only discrete valuation. Proposition 2.11, p8:L4, is an \textbf{envelope-only finite normal ground-field extension} statement; it is not descent of unbounded operators along an alteration of varieties.

\noindent \textbf{[P]} Pille--Schneider, arXiv:2209.04879v2, Remark 4.4, p35:L36, distinguishes defining finite energy over any field from extending its operator in general. It is used only as a scope cross-check. The remaining four supplied caches and the trivial-valuation [T] cache have their hashes checked; they supply no indispensable new argument here. Thuillier is a curve theory; the Berkovich 1993 cache is the etale-cohomology paper; the CLD and GR caches are not an imported general BFJ finite-energy surjectivity theorem. No book or other campaign was inspected.

\section*{3. What perturbation-friendly ideals actually replace}

We check the core mixed-characteristic envelope proof on the original variety. Suppose \(\mathscr X\) is a projective regular \(R\)-model and \(\mathscr L\) extends \(L\). Assume \(L\) also has an ample extension \(\mathscr A\) on this same model. Positive powers allow integral notation. Let \(D\) represent \(\mathscr L\), and put

\[
a_m=\operatorname{Im}\bigl(H^0(\mathscr X,\mathscr L^m)\otimes\mathscr L^{-m}\to\mathcal O_{\mathscr X}\bigr),
\quad b_m=\tau(\mathcal O_{\mathscr X},m\Vert D\Vert),
\quad t_0=\tau(\mathcal O_{\mathscr X},0),\quad a'_m=a_mt_0.
\]

For all sufficiently large m the generic base ideal is the unit ideal, because L is ample; hence a\_m is vertical. By [F], Remark 4.2(v), t\_0 is vertical, \textbf{not assumed to be the unit ideal} even on a regular mixed-characteristic model. Coherence and quasi-compactness imply \(\pi^N\mathcal O\subset t_0\subset\mathcal O\) for some N. Thus \(g_0=\log|t_0|\) is bounded, between \(N\log|\pi|\) and zero. This is the precise harmless correction, not an assertion that mixed-characteristic test ideals behave identically to multiplier ideals.

By [F], Lemmas 4.4 and 4.6, and Theorem 4.7(ii),

\[
a'_m\subset b_m,\qquad b_{m\ell}\subset b_m^\ell.
\]

These are proved perturbation-friendly subadditivity statements; Conjecture 3.6 about older +-ideals is unnecessary. To check uniform generation, choose a globally generated ample divisor H and then choose \(m_0\) such that

\[
\mathscr A^{m_0}\otimes\mathcal O(-K_{\mathscr X}-(n+1)H)
\quad\hbox{is globally generated}.
\]

In Theorem 4.7(i), take \(\lambda=m\) and \(E=mD+H+K_{\mathscr X}\). Its required difference is

\[
E-K_{\mathscr X}-\lambda D=H,
\]

which is big and nef. Here n is the \textbf{special-fiber dimension}, not the total dimension of an equal-characteristic spread-out family. The theorem and the preceding generated twist give global generation of \(b_m\mathscr A^{m_0}\mathscr L^m\), independently of m. This verifies the actual condition in the body of [F], p9, rather than importing the different equal-characteristic condition from [G].

After twisting \(\mathscr A\) by a principal multiple of the fiber, arrange \(\mathscr A=\mathscr L(-E_0)\) with \(E_0\) effective vertical. Let \(g_{E_0}\) be the bounded model function of \(E_0\), with \(\log|\mathcal O(-E_0)|=-g_{E_0}\). Global generation implies

\[
a_{m+m_0}\supset b_m\mathcal O(-m_0E_0).
\]

Write \(\phi_m=m^{-1}\log|a_m|\). The pointwise limit \(\phi_m\to P_\theta(0)\) follows from the discrete model-envelope base-ideal argument [G], Proposition 2.10. The correction \(g_0/m\) tends uniformly to zero. Taking \(\ell\to\infty\) in \(a'_{m\ell}\subset b_m^\ell\) therefore gives

\[
m^{-1}\log|b_m|\geq P_\theta(0).
\]

Combining these inequalities, for large k with \(k-m_0>0\), yields

\[
\phi_k\leq P_\theta(0)
\leq\frac{k}{k-m_0}\phi_k+
\frac{m_0}{k-m_0}g_{E_0}.
\]

The large-index base ideals are vertical and their log functions are finite continuous functions. Superadditivity \(a_ma_\ell\subset a_{m+\ell}\), using finitely many large starting indices to cover every residue class modulo a fixed large index, gives a fixed lower bound \(\phi_k\geq-C\). Since \(\phi_k\leq0\),

\[
0\leq P_\theta(0)-\phi_k
\leq\frac{m_0}{k-m_0}(C+\Vert g_{E_0}\Vert_\infty)\longrightarrow0.
\]

This proves the uniform approximation on this model. Only the large-index tail is used; no assertion about vertical base ideals for small non-generated powers is necessary. Theorem 7.1 extends it to every original continuous direction via uniform approximation by model functions, envelope translation, and Lemma 6.4's regular ample extension. Lemma 6.4 retains projective resolution. When adjusting its relatively ample twist to preserve the generic line bundle, choose the subsequent power of the ample base bundle sufficiently large after that adjustment; relative ampleness is unchanged by a pullback twist.

Thus [F] replaces the ideal-theoretic input. It has not supplied regular projective domination or embedded resolution. The hypothesis that one good model exists is insufficient for its all-directions application on all dominating model determinations.

\section*{4. An unconditional model-level theorem in the whole retained domain}

\textbf{Proposition 1 (test orthogonality and target approximation).} For every original \((K,X,L)\) of Section 1, with no residue-characteristic restriction, the synthetic test datum has orthogonality. For every positive Radon \(\mu\) of mass V and finite \(\mathcal E^*_\theta\), there is a sequence \(v_j\in D_\theta\) such that the genuine raw model measures \(\MA_\theta(v_j)\) converge weakly to \(\mu\), and their normalized dual energies converge to \(A\). This does not assert that an unbounded psh potential realizes the limit.

\textbf{Proof.} [M], Sections 1.2.2 and 1.3.2, provide the datum on the actual compact \(X^{\rm an}\): D is dense in \(C^0\), contains constants, and the form order is nice. On a finite-dimensional set of forms, a common model expresses nefness by closed curve-degree inequalities; positive domination makes the positive cone span. The intersection pairing is symmetric, positive on positive tuples, and its vertical integration-by-parts form is seminegative. The latter is the characteristic-unrestricted local Hodge theorem imported in [M], not an application of SNC resolution here. The positive generic ample class lies in \(\operatorname{Pos}(X)\) as stated there. Geometric integrality implies the required irreducibility for the submean property whenever it is used. These check the source's hypotheses, without presuming a general cohomology identification.

Real psh model tests can be approximated uniformly by rational ones: on a common higher model choose an ample representative of the same generic ample class, perturb toward it, approximate the finitely many vertical coefficients rationally inside the ample cone, and let the perturbation decrease. Strict constant margins and rational model maxima then give the precise directed-maxima condition [M], Definition 2.11. [M], Example 2.20, applies to this projective Berkovich datum and supplies

\[
\lim_{v\in D_\theta,\ v<f}\int(f-v)m(v)=0\qquad(f\in D).
\]

This is the imported published test-level conclusion. It does not assert continuity or psh membership of the pointwise envelope on an arbitrary model system.

For completeness the passage from this input to weak approximation is explicit. Define

\[
a_\nu(v)=A-e(v)+\int v\,d\nu\geq0.
\]

Choose maximizing tests \(v_j\), so \(a_\nu(v_j)\to0\), directly from the finite supremum. For a fixed \(f\in D\), [M], Proposition 2.19, gives a constant \(C_f\) independent of v and t with

\[
\left|\widetilde e(v+tf)-e(v)-t\int f\,m(v)\right|\leq C_ft^2,
\]

where \(\widetilde e(h)=\sup\{e(w):w\in D_\theta,\ w\leq h\}\). The dual definition gives \(\widetilde e(v_j+tf)\leq A+\int(v_j+tf)d\nu\). Apply these inequalities for positive t and then for -f to obtain

\[
\left|\int f(m(v_j)-\nu)\right|
\leq a_\nu(v_j)/t+C_ft\qquad(t>0).
\]

Let j tend to infinity and then t tend to zero. Density of D in \(C^0\) and probability mass give weak convergence. Uniform concavity at the midpoint, [M], Theorem 1.33, gives, with

\[
J(u,v)=e(u)-e(v)+\int(v-u)m(u),\qquad
J(v)=\int v\,m(0)-e(v),
\]

the bounds \(J(u,v)\leq C_n(a_\nu(u)+a_\nu(v))\) and \(J(v)\leq C_n(A+a_\nu(v))\). The target's linear terms cancel in that concavity inequality. Hence the maximizing sequence has bounded J and is Cauchy for this test quasi-distance. The mixed Holder estimate of [M], used in Theorem 2.23, implies \(\int v_j(m(v_j)-\nu)\to0\); this is exactly the additional argument there. Therefore

\[
e^*(m(v_j))=e(v_j)-\int v_jm(v_j)\longrightarrow A.
\]

Multiplication by V proves the raw assertion. All measures in this conclusion are intersection-theoretic model Monge--Ampere measures. The proof constructs an approximating sequence of such measures, not an inverse measure by definition. It ends here in the full original domain: no normalized compactness of the full potential cone, or an operator on its bounded discontinuous members, was supplied. In particular neither the preceding estimate nor completeness of the measure space justifies an actual E1 solution. \(\square\)

\section*{5. A genuine mixed-characteristic equation theorem with explicit conditional status}

\textbf{Theorem 2 (conditional partial extension).} Let K be any complete nontrivially discretely valued mixed-characteristic field, and retain arbitrary-dimensional smooth geometrically integral projective X and the actual ample rational L. Additionally assume \textbf{both} [F], Definition 6.1 and Definition 6.3 for the projective models of this particular X. Then the full SNC-based psh cone has a genuine approximation-independent bounded and finite-energy mixed local product theory. The diagonal finite-energy operator has raw mass V. Its image is exactly the positive Radon measures of mass V and finite dual energy in Section 1. For each such measure, with arbitrary support, there is an E1 solution, unique modulo additive real constants. No continuity of the unbounded solution is asserted.

This is a conditional theorem, not a solution of original P16.3. Its hypotheses quantify over every projective model of X. They cannot be replaced by existence of one regular model, locally projective arithmetic resolution, or a supplied alteration of the generic variety.

\textbf{The full cone and compactness.} Under the stated assumptions every projective model is dominated by a projective SNC model, by [F], Section 8. On every such determination of theta, let \(p_{\mathscr X}\) be its retraction to its dual complex. Define \(\PSH(X,\theta)\) to be the usc functions u such that \(u\leq u\circ p_{\mathscr X}\), and such that u restricted to each dual complex is a uniform limit of restrictions of theta-psh model functions, exactly [S], Definition 7.1. These restrictions are finite and continuous; the cone contains all model psh potentials and is stable under convex combinations, max, constants, and decreasing limits. This is the full cone on X, not a selection of functions factoring through one skeleton.

Its weak topology is uniform convergence on each dual complex, equivalently pointwise on divisorial points in this cone. [F], Section 8, transfers [S], Sections 1--7, in particular normalized compactness, continuity of sup, and determination by divisorial values. Integration against any fixed model MA measure is continuous since its support is a finite set of such points. Section 3 supplies all continuous envelopes and uniform model approximation. The same envelope and skeleton argument as [G], Corollary 9.1, or the explicit import of [S], Theorem 8.7 in [F], gives decreasing model regularization. The [B]-v1 sequence construction also applies; arbitrary decreasing nets are retained for continuity and compactness arguments.

\textbf{Bounded operator.} On model tuples use Section 1's mixed intersections. Under this regularization, the induction in [B], Theorem 3.1, constructs the unique positive mixed measure for bounded psh tuples, agrees with model intersections, and gives decreasing-net continuity with uniformly bounded continuous test weights. The proof uses the model projection and integration-by-parts identities and seminegative vertical pairing; it uses no multiplier or test ideal. Bounded locality is [B], Theorem 5.1, whose proof uses that bounded construction, max and regularization. Thus these results transfer with their exact hypotheses now supplied on X. Uniform approximation of continuous psh metrics agrees with the Chambert--Loir operator; the extended product is genuine. Raw mixed mass is V on every bounded tuple, by decreasing continuity and testing the constant function 1.

Extend e by the infimum of the energies of dominating model psh potentials. [B], Proposition 6.2, using the same regularization and intersection pairing, yields increasing concave weakly upper semicontinuous energy, continuous along decreasing nets, and \(e(u+c)=e(u)+c\). Set

\[
\mathcal E^1(X,\theta)=\{u\in\PSH(X,\theta):e(u)>-\infty\}.
\]

No completeness of a space of measures is used in this definition or to construct the following operator.

\section*{6. Unbounded products, mass, integrability, and weighted limits in Theorem 2}

Normalize u to be nonpositive, put \(u_t=\max(u,-t)\), and define

\[
\MA_\theta(u)=\lim_{t\to\infty}
\mathbf1_{\{u>-t\}}\MA_\theta(u_t),\qquad
m(u)=V^{-1}\MA_\theta(u).
\]

Bounded locality gives compatibility on retained sets, so these are increasing positive restricted measures. Their limit is Radon, is independent of constant normalization, and charges no pluripolar set. The pole-set assertion follows from the bounded calculus [B], Proposition 3.7 and Remark 6.5, not from an arbitrary weak-limit claim.

The bounded energy difference formula gives

\[
t\,m(u_t)\{u\leq-t\}
\leq2(n+1)\bigl(e(u_{t/2})-e(u_t)\bigr)\longrightarrow0.
\]

Indeed \(u_{t/2}-u_t=t/2\) on that set and its integral against \(m(u_t)\), divided by n+1, is one of the nonnegative terms of the difference formula. Decreasing energy continuity proves the limit. Thus the restricted measures have limiting raw mass \textbf{exactly} V. For continuous f, comparison of two truncations and \(\int(-u_t)m(u_t)\leq(n+1)|e(u)|\) give

\[
\left|\int f\,m(u_t)-\int f\,m(u)\right|
\leq\frac{2(n+1)|e(u)|}{t}\Vert f\Vert_\infty.
\]

These bounds are uniform for a decreasing net \(u_j\downarrow u\in\mathcal E^1\), since \(e(u_j)\geq e(u)\), after restricting to a bounded-above tail of the net. Bounded decreasing continuity then proves \(m(u_j)\to m(u)\). Approximation independence concerns every decreasing approximation, not just a selected model sequence.

For n possibly different inputs \(u_i\in\mathcal E^1\), simultaneous bounded truncations give mixed measures \(T_t\). Let \(h_t=n^{-1}\sum_i u_{i,t}\) and \(h=n^{-1}\sum_i u_i\). The positive expansion of \(\MA(h_t)\) gives

\[
T_t\leq\frac{n^n}{n!}\MA(h_t).
\]

On \(B_t=\bigcup_i\{u_i\leq-t\}\), one has \(h_t\leq-t/n\). Concavity bounds its energy below by finite \(e(h)\), hence

\[
T_t(B_t)\leq\frac{n^n}{n!}\frac{n(n+1)V|e(h)|}{t}.
\]

On the complement bounded mixed locality gives compatible restricted products. Their increasing limit has mass V. It equals the weakly polarized limit: diagonals of all convex combinations have already proved decreasing continuity, and finite polynomial interpolation determines every mixed coefficient. This proves positivity and independence, and identifies the product as the actual restricted local product rather than only an abstract polarization.

For u,w nonpositive in E1, the proof of [B], Lemma 6.10, gives the correctly signed bound

\[
\int(-w)m(u)
\leq 2^{n+1}(n+1)\bigl[-e((u+w)/2)\bigr]<\infty.
\]

It follows by the bounded mixed energy inequality with w\_s, first decreasing the weight and then increasing the retained u-truncation region. The displayed sign in the v1 lemma is not used. Applying this to a convex average of the inputs and their positive mixed expansion proves integrability of every E1 weight against every mixed E1 measure.

Here is an explicit weighted-limit verification that avoids testing weak convergence against an unbounded weight. First take decreasing model inputs with mixed measures T\_j and averages h\_j. Their suprema have a common upper bound and their energies are bounded below by e(h). The same holds for a decreasing model weight sequence w\_l with limit w in E1. Consequently all of their model J values have a common finite bound R. For k greater than or equal to l, concavity gives

\[
0\leq J(w_l,w_k)\leq e(w_l)-e(w_k)\leq e(w_l)-e(w)\longrightarrow0.
\]

The model-level mixed Holder estimate [M], equation (1.37), applied to \(w_l,w_k\) and \(h_j,0\), therefore gives a constant independent of j,l,k such that, with \(\alpha=2^{-n}\),

\[
0\leq\int(w_l-w_k)m(h_j)
\leq\int(w_l-w_k)m_0+
C_R\bigl(e(w_l)-e(w)\bigr)^{2^{-n}}.
\]

Here the nonnegativity comes from \(w_l\geq w_k\). Model measures have finite divisorial support, and w is finite on that support. Letting k increase to infinity by monotone convergence is thus legitimate. The right side tends to zero as l increases, uniformly in j: decreasing energy continuity handles the second term and monotone convergence against the fixed \(m_0\) handles the first. The positive expansion bound \(T_j\leq(n^n/n!)\MA(h_j)\) proves

\[
\lim_{l\to\infty}\sup_j\int(w_l-w)T_j=0.
\]

For fixed l, w\_l is continuous, so weak mixed continuity gives convergence of its integrals. Passing first j and then the fixed continuous weight index k in the displayed estimate gives the same tail bound for the limiting mixed measure T. Monotone convergence proves w is integrable against T and yields convergence of the integrals of w against T\_j. If the weight decreases simultaneously with the inputs, its difference from w is bounded by the same uniform tail. This proves the joint weighted assertion. For general decreasing E1 inputs, apply the same model estimates on simultaneous dominating model approximations; their energies and m\_0 integrals converge. The uniform bound just proved survives that passage and the already established weak continuity identifies their measures. Directed indices and eventual common upper bounds give the identical argument for decreasing nets. Only E1 psh weights and their real span are asserted, not arbitrary unbounded Borel weights.

Integration by parts and Cauchy--Schwarz now pass from models to E1 using these actual weighted limits. Equivalently, for a finite tuple the energies of convex combinations are polynomials of degree at most n+1; decreasing energy continuity and finite interpolation give coefficient convergence. The weighted-limit proof identifies the coefficients with genuine mixed integrals, so there is no separate abstract mixed product at this step. This supplies the finite-energy calculus used in [B], Proposition 8.2, with its prerequisites verified here.

Bounded locality on the retained truncation sets, full mass, and the zero pole mass give

\[
\mathbf1_{\{u>v\}}m(\max(u,v))=\mathbf1_{\{u>v\}}m(u),
\qquad\int_{\{u<v\}}m(v)\leq\int_{\{u<v\}}m(u).
\]

The comparison passage uses \(\max(u,v-\epsilon)\) and then \(\epsilon\downarrow0\), so equality-locus mass is not discarded without justification. The mixed polynomial formula and weighted continuity yield

\[
\left.\frac{d}{dt}\right|_{0+}e((1-t)u+tv)=\int(v-u)m(u).
\]

\section*{7. Actual envelope orthogonality and every continuous direction}

Here Proposition 1 supplies a replacement even for the orthogonality proof used in [F], so the separately cited uncached [BGJ+20] is unnecessary as an additional application. For \(f\in D\), the increasing directed net of tests v<f has supremum \(P_\theta f\), which is continuous by Section 3 under Theorem 2's assumptions. Dini convergence for nets is uniform. Uniform continuity of the genuine continuous-metric MA operator and Proposition 1 therefore give

\[
\int(f-P_\theta f)m(P_\theta f)=0.
\]

Uniform approximation of arbitrary \(f\in C^0\) by D and the envelope Lipschitz inequality extend this identity to every original continuous f. Thus [B], Theorem 7.2, applies with its characteristic-zero orthogonality proof replaced.

For any u in E1 and \(f\in C^0(X^{\rm an})\),

\[
u-|t|\Vert f\Vert_\infty\leq P_\theta(u+tf)
\leq u+|t|\Vert f\Vert_\infty.
\]

The perturbed envelope is in E1. A decreasing model approximation \(u_j\downarrow u\) gives decreasing envelopes \(P(u_j+tf)\downarrow P(u+tf)\): the limit is psh, below u+tf, and dominates its envelope. The integrated continuous-envelope identity, bounded by \(\Vert f\Vert_\infty\), passes to the limit by decreasing operator/energy continuity and dominated convergence in its real parameter. It yields

\[
e(P(u+tf))-e(u)=\int_0^t\left(\int f\,m(P(u+sf))\right)ds.
\]

For positive s, replace f by \(f+\Vert f\Vert_\infty\). The resulting envelopes decrease to u as s decreases to zero; their operators therefore converge to m(u). Subtracting the constant s-shift does not change MA. For negative s use \(f-\Vert f\Vert_\infty\) and the same monotonicity argument. Consequently the derivative is genuinely two-sided:

\[
\left.\frac{d}{dt}\right|_0 e(P_\theta(u+tf))=\int f\,m(u)
\qquad\text{for every original }f\in C^0(X^{\rm an}).
\]

No target support or a single skeleton is used in this passage.

\section*{8. Finite dual energy, genuine attainment, and arbitrary support}

For a target satisfying Section 1, every u in E1 is integrable against \(\nu\). Indeed, choose decreasing model tests \(v_j\downarrow u\). They are uniformly bounded above, while

\[
\int v_j\,d\nu\geq e(v_j)-A\geq e(u)-A.
\]

Monotone convergence proves integrability and

\[
e(u)-\int u\,d\nu\leq A.
\]

It also proves that the dual supremum over D equals the supremum over all E1. Conversely, mixed integrability and the tangent inequality give

\[
e^*(m(u))=e(u)-\int u\,m(u)<\infty.
\]

The inequality follows from energy concavity; equality follows by decreasing model approximation to u. Thus the claimed image and target conditions coincide intrinsically.

We prove attainment in the actual potential space. Put \(m_0=m(0)\). Compactness and continuity of its finite-divisorial-support integral give a finite T with

\[
-T\leq\int u\,m_0\leq0\quad(\sup u=0).
\]

For C>0 the sets

\[
K_C=\{u\in\PSH(X,\theta):\sup u=0,\ e(u)\geq-C\}
\]

are weakly compact and lie in E1, with \(0\leq J(u)\leq C\). Choose maximizing model tests r\_j. Proposition 1's defect estimate bounds their J. [M], Lemma 2.25, with the other test zero, gives, for a fixed common energy bound R,

\[
\left|\int v\,(\nu-m(r_j))\right|
\leq C_{n,R}\,a_\nu(r_j)^{1/2}
\]

for tests with \(J(v)\leq R\). To extend it to u in K\_C, take decreasing model \(u_\ell\downarrow u\). Both \(e(u_\ell)\) and the m\_0 integral converge, so eventually \(J(u_\ell)\leq C+1\). Monotone convergence in both integrals is legitimate: the target integrability was just proved and m(r\_j) has finite divisorial support. Although the eventual index can depend on u, the bound does not. Hence

\[
\sup_{u\in K_C}\left|\int u\,d\nu-\int u\,m(r_j)\right|
\leq C_{n,R}a_\nu(r_j)^{1/2}\longrightarrow0.
\]

The approximating functionals are weakly continuous on K\_C. Their uniform limit proves continuity of target integration there. This is the indispensable potential-space passage; the measure completion theorem does not imply it without the supplied compactness and regularization.

Normalize a maximizing sequence v\_j by sup=0. Its J is bounded and

\[
-e(v_j)=J(v_j)-\int v_jm_0\leq J(v_j)+T.
\]

It lies in one K\_C. Compactness gives a convergent \textbf{subnet}, with limit u in K\_C. Upper semicontinuity of energy and continuity of target integration yield

\[
e(u)-\int u\,d\nu\geq\limsup_j\left(e(v_j)-\int v_j\,d\nu\right)=A.
\]

The reverse inequality was proved above. Thus u is an actual E1 maximizer on X. For arbitrary original continuous f set \(u_t=P(u+tf)\) and

\[
h(t)=e(u_t)-\int(u+tf)\,d\nu.
\]

Since \(u_t\leq u+tf\), positivity and maximality give \(h(t)\leq A=h(0)\). Section 7's two-sided derivative gives

\[
0=h'(0)=\int f\,m(u)-\int f\,d\nu.
\]

Testing all continuous f proves \(m(u)=\nu\), hence \(\MA_\theta(u)=\mu\) with raw mass V and unrestricted support. No artificial inverse operator or replacement measure was defined.

\section*{9. Uniqueness and a characteristic-independent strict kernel}

Let u,v in E1 have identical diagonal operator, put f=u-v, and set \(S_0=\int f(m(v)-m(u))=0\). All integrals below are finite by Section 6. For a semipositive model reference \(\eta=\theta+dd^cg\) in the same actual ample class, curvature identities on bounded approximants and decreasing continuity give the same genuine operators for u-g and v-g. The bounded energy difference formula and decreasing continuity give the exact cocycle

\[
e_\eta(u-g)=e_\theta(u)-e_\theta(g).
\]

Thus this is a justified change of reference on X, not descent through another variety. Write \(\eta_u=\eta+dd^c(u-g)\), \(\eta_v=\eta+dd^c(v-g)\), and for \(i+j+k=n-1\) put

\[
q_{i,j,k}=-\int f\,dd^cf\wedge\eta_u^i\wedge\eta_v^j\wedge\eta^k\geq0.
\]

For k=0 these norms sum to \(VS_0=0\), by telescoping the diagonal products. The field-independent induction of [B], Proposition 8.2, now gives all q=0. At zero its key step can be stated directly: expand \(\eta^k-\eta_u^k\) as \(-dd^c(u-g)\sum_{\ell=0}^{k-1}\eta_u^\ell\eta^{k-1-\ell}\), multiply by the remaining factors, and integrate. In the resulting cross term replace dd\textasciicircum{}cf by eta\_u minus eta\_v. Cauchy--Schwarz pairs f and u-g against each positive product containing one of eta\_u,eta\_v and fewer than k eta factors. Its f-square norm is a previously vanishing q, and the other square norm is finite by the verified mixed calculus. Thus every cross term is zero. This proves \(q_{0,0,n-1}=0\). For general S\_0 the same induction has the bound \(C S_0^{2^{1-n}}\), with V absorbed in C; only the zero conclusion is needed.

On any sufficiently high regular SNC model choose eta represented by a relatively ample rational class A with generic fiber L, using [F], Lemma 6.4 and Remark 6.5 and a vertical reference change. Write the special fiber as \(\sum_i b_iE_i\). For a vertical real divisor \(D=\sum_i a_iE_i\), the principal-fiber relation gives

\[
D^2A^{n-1}=-\sum_{i<j}b_ib_jw_{ij}
 \left(\frac{a_i}{b_i}-\frac{a_j}{b_j}\right)^2,
\qquad w_{ij}=(E_iE_jA^{n-1}).
\]

Indeed expand the square and use \((\sum_i b_iE_i)E_jA^{n-1}=0\) to eliminate diagonal terms. The effective intersections and ampleness make w\_ij nonnegative, and strictly positive when the two components meet. This calculation works in every characteristic and for n=1 as well. Mixed-nef factors instead of A give nonnegative weights and seminegativity, without the strictness assertion.

The graph is connected. The generic fiber is proper geometrically integral, hence its global functions are K. Properness makes the global-functions algebra of the normal integral model a finite R-module inside K, so it is integral over R; since R is integrally closed it equals R. Stein factorization gives connected special fiber. Consequently zero of the displayed form forces \(a_i/b_i\) constant over every component.

Let \(\lambda=-\log|\pi|\) and choose \(a_i=b_if(x_i)/\lambda\). The associated model function g\_D equals f at the component vertices. Cauchy--Schwarz and \(q_{0,0,n-1}=0\) give

\[
\int f\,dd^cg_D\wedge\eta^{n-1}=0
=\lambda D^2A^{n-1}.
\]

The second equality is the actual model-intersection identity with the indicated metric convention. Hence f has one constant value on those vertices. Keep a fixed divisorial anchor while refining models. Cofinality places any other divisorial point with that anchor on a dominating SNC model, aligning the constants even though the ample reference varies. Divisorial determination gives u=v+c on all X. If both suprema are zero, c=0. A common model-function subtraction is not claimed to preserve equality of normalized suprema. The metrics differ by the positive factor exp(-c). This completes the proof of the \textbf{conditional} Theorem 2.

\section*{10. A concrete finite-cover averaging failure}

This example is a diagnostic mathematical obstruction to a proposed proof strategy, not an endpoint substituted for the general variety in P16.3.

Take \(K=\mathbb Q_5\), \(R=\mathbb Z_5\), and

\[
Y=\mathbb P^1_K\times\mathbb P^1_K,
\quad X=\mathbb P^1_K\times\mathbb P^1_K,
\quad q(z,w)=(z^2,w),\quad\sigma(z,w)=(-z,w).
\]

The homogeneous square map is a finite flat degree-two map. Let L=O(1,1) on X and \(H=q^*L=O(2,1)\) on Y, of raw volume 4. On the good projective R-model of Y blow up the two distinct special-fiber points \((z,w)=(1,0),(-1,0)\). The blowup is projective regular and its special fiber has multiplicity-one SNC components: the strict transform S and two exceptional planes E\_1,E\_2. The involution swaps the two planes.

Let \(\epsilon=1/2\), and use the rational model classes

\[
A_1=H-\epsilon E_1,\qquad A_2=H-\epsilon E_2,
\qquad A_w=H-\tfrac\epsilon2(E_1+E_2).
\]

They are relatively nef. On each exceptional plane the restriction of H is trivial, \(\mathcal O(E_i)|_{E_i}=O(-1)\), and the other exceptional divisor is disjoint, giving a nonnegative O(1) multiple. On S, the restricted divisors are the corresponding pullbacks of O(2,1) minus the indicated exceptional curves. To check nefness directly, use the special-fiber curve F defined by w=0. For any integral curve C different from F, the sum of its multiplicities at the two centers is at most \(F\cdot C\leq H\cdot C\); each coefficient subtracted is at most 1/2. For F itself, \(H\cdot F=2\), and at most 1 is subtracted. Exceptional curves on S have degree either 0 or a positive coefficient. These checks cover all special-fiber curves, proving nefness without an assumed resolution theorem.

Let u and \(\sigma^*u\) be the actual model potentials of A\_1 and A\_2 relative to H, and let \(w=(u+\sigma^*u)/2\) be the potential of A\_w. The raw products follow by restricting to components. On S, exceptional curves have square -1, are disjoint, and have zero intersection with the pulled-back H. On E\_i, the relevant square has degree \(\epsilon^2\). Therefore

\[
\MA_H(u)=(4-\epsilon^2)\delta_S+\epsilon^2\delta_{E_1},
\quad
\MA_H(\sigma^*u)=(4-\epsilon^2)\delta_S+\epsilon^2\delta_{E_2},
\]

and the mixed product of A\_1,A\_2 is \(4\delta_S\). Thus

\[
\MA_H(w)=(4-\epsilon^2/2)\delta_S+
 \frac{\epsilon^2}{4}(\delta_{E_1}+\delta_{E_2}),
\]

whereas the averaged target is

\[
\tfrac12(\MA_H(u)+\MA_H(\sigma^*u))=
(4-\epsilon^2)\delta_S+\frac{\epsilon^2}{2}(\delta_{E_1}+\delta_{E_2}).
\]

At epsilon=1/2 the two vectors, ordered as S,E\_1,E\_2, are respectively

\[
(31/8,1/16,1/16),\qquad (15/4,1/8,1/8),
\]

with difference \((1/8,-1/16,-1/16)\). Both measures have raw mass 4. Hence averaging genuine semipositive model metrics does \textbf{not} average their MA measures, even in this explicitly regular example. All potentials here are bounded model potentials and consequently finite energy; no unbounded descent issue causes the failure.

The conclusion is limited: it disproves arbitrary-solution averaging as a way to solve an averaged target. It does not disprove descent of an already invariant solution, or descent obtained by pulling up an invariant target and using a proved uniqueness theorem. Those routes still need the actual pullback target, descent of the metric cone, cocycle, energy, and operator identities. The finite arithmetic check in descent\_check.py confirms the displayed fractions, not the geometric proof.

\section*{11. Covers, inseparability, and the missing descent package}

For a finite surjective q of degree d between integral projective n-folds, the projection formula on compatible models with \textbf{pulled-back metrics} gives the valid raw identities

\[
(q^*L)^n=dV,\qquad
q_*\MA_{q^*\theta}(q^*v_1,\ldots,q^*v_n)=d\MA_\theta(v_1,\ldots,v_n),
\]

and, for model potentials,

\[
\mathcal E_{q^*\theta}(q^*v)=d\mathcal E_\theta(v),
\qquad e_{q^*\theta}(q^*v)=e_\theta(v).
\]

Indeed each weighted intersection is the pullback of the original metric data and pushes forward to d times the corresponding intersection cycle; summing the energy formula gives the last identities. These are forward identities for genuine model metrics, not an unproved reverse descent statement for arbitrary E1 data. If both sides already have compatible decreasing regularization and finite-energy calculus, they extend to pulled-back E1 potentials by that regularization. The unrestricted original case has not supplied those prerequisites.

A radicial map does not remove the degree factor. In equal characteristic p, the relative coordinate Frobenius on \(\mathbb P^n_K\) is finite radicial of degree \(p^n\), sends O(1) to O(p), and induces a universal homeomorphism on the underlying spaces. In the analytic statement, a valuation extends uniquely across a purely inseparable residue-field extension: the value of an element whose p-th power is in the smaller field must be that power's value to the exponent 1/p. Existence and uniqueness give a continuous bijection; proper analytic spaces are compact Hausdorff, so it is a homeomorphism, also after base change. Its pulled-back model MA measure has raw mass \(p^n\), while the original mass is 1. The model identities above still have d=p\textasciicircum{}n. The topological identification is not a proof that genuine metrics, normalized energies, raw measures and their local products descend with degree one.

An alteration can create an additional polarization problem. If a proper generically finite morphism q contracts a curve C, then \((q^*L)\cdot C=0\), so q\textasciicircum{}*L is not ample. Applying an ample equation theorem on the altered variety with q\textasciicircum{}*L is then unjustified. Adding a small ample class changes the curvature and volume and requires a separately proved limiting equation and descent theory in the resulting nonample class. This document does not import an unproved P16.5 theorem to close P16.3.

Even for a finite Galois cover preserving ampleness, a successful route must supply: regular or replacement model theory on the cover \textbf{for every relevant model}; a genuine pullback measure with the correct local degrees and finite dual energy against all cover tests, not just pullback tests; uniqueness on the cover forcing invariance; descent of the invariant singular metric as a psh metric on the original L; compatible model approximations; raw mass and mixed local products after descent; the energy cocycle and finite-energy integrability; every original C0-direction derivative and comparison. The model energy identity alone controls only pullback tests, a subset of all tests on the cover. An envelope commutation statement for a finite normal ground-field extension in [G] does not establish this package for a generically finite alteration of varieties.

No such all-model cover package is constructed here. The averaging counterexample identifies a real failed passage, but supplies no impossibility theorem for all conceivable cover approaches.

\section*{12. Concrete obligations still belonging to original P16.3}

The unchanged general claim remains unproved in this work. The following gaps are indispensable rather than technical provenance issues.

\noindent 1. \textbf{Mixed characteristic in arbitrary dimension.} [F] now supplies the ideal-theoretic replacement and continuous envelopes once its model-resolution premise is available. This work does not prove [F], Definition 6.1 for every original X, nor Definition 6.3, nor a replacement that gives the full cone, compactness, bounded locality, unbounded products and comparison on X without them. The locally projective arithmetic-threefold statement cited in [F], Remark 6.2 is not a checked projective domination theorem. Theorem 2 therefore remains conditional even for general mixed-characteristic surfaces.

\noindent 2. \textbf{Equal characteristic in arbitrary dimension.} The existing [G] geometric-origin range retains both resolution and embedded resolution in total dimension d+n; only its specified low-dimensional range is supplied unconditionally. No arbitrary-dimensional projective model replacement is proved here. A test-level orthogonality theorem does not supply SNC cofinality or a bounded discontinuous product.

\noindent 3. \textbf{Equal-characteristic fields and varieties outside geometric origin.} No realization of every original X over an allowed family is proved. No replacement envelope continuity, potential compactness, or full unbounded calculus is constructed there. The mixed-characteristic regular-model test ideals do not transfer to this regime just by changing the characteristic adjective.

\noindent 4. \textbf{Potential-space realization on the full domain.} Proposition 1 gives authentic approximating model measures for all finite-dual-energy targets. To complete the original task one must prove a full psh cone with an adequate normalized compactness theorem, compatible bounded local products and a finite-energy extension with raw mass V, arbitrary decreasing and integrable weighted limits, and the two-sided envelope derivative in every original C0 direction. Only then can the actual attainment and uniqueness proof of Sections 8--9 be applied without conditional geometry. Declaring a measure completion to be a space of potentials does not prove these facts.

\noindent 5. \textbf{Any proposed finite/altered descent.} All reverse descent requirements in Section 11 remain unproved in the general domain. Nonlinearity, degree factors under inseparability, and loss of ampleness under a nonfinite alteration must be handled in an actual proof, not ignored.

The new source audit removes the obsolete claim that perturbation-friendly subadditivity itself is missing. It also isolates a test-level orthogonality input that already holds without resolution. It does not prove the remaining geometry and potential realization, and it does not prove that those gaps are impossible to close. The retained scope is not weakened to count Theorem 2 or Proposition 1 as a complete answer.

\section*{13. Reproducibility and evidential limits}

The frozen input SHA256 is \texttt{3de21500\allowbreak{}70e5541d\allowbreak{}b27997d1\allowbreak{}bd276df3\allowbreak{}c0604536\allowbreak{}783ec3ed\allowbreak{}ade44b37\allowbreak{}a2dcce8c}. It differs from the historical artifact hash because this is a new, larger frozen assignment. input.json and input-original.json are byte-identical; assignment-original.txt is the decoded assignment string; task-original.json preserves the decoded task object. The historical proof and review can be recovered from the decoded assignment without reading the prior workspace. No original public cache is modified.

Rerun from this workspace:

\begin{verbatim}
python3 audit_sources.py
python3 descent_check.py
python3 build_paper.py
python3 check_artifacts.py
\end{verbatim}

audit\_sources.py verifies the exact twelve authorized sources and retains literal page extracts, PDF page/line anchors, hashes and versions in source-extracts/source-provenance.json. source-application-ledger.json records the precise applications and unresolved hypotheses. descent\_check.py uses exact rational arithmetic for Section 10; its observations have no force beyond the displayed coefficient check. build\_paper.py generates the standalone P16.3-paper.tex from this argument and executes two actual TeX passes, retaining both stdout/stderr streams, individual TeX logs, the PDF, and command return codes. environment.json records actual installed tool versions; no software is installed or runtime setting changed. check\_artifacts.py records input preservation, schema, source and build checks and an artifact manifest.

The mathematical evidence is the cited literal primary-source conclusions and the deductions with explicit hypotheses above. Hash verification, exact arithmetic, file checks and successful compilation are finite observations about artifacts. They do not establish mathematical correctness or close original P16.3. No other model, subagent, external message, network retrieval or heavy computation was used. status.json has exactly the required five fields and status=partial.

\end{document}
